\documentclass[11pt]{article}
\usepackage[margin=1in]{geometry}
\usepackage{amsmath,amssymb,amsthm}
\usepackage[hidelinks]{hyperref}
\newtheorem{theorem}{Theorem}
\newtheorem{lemma}[theorem]{Lemma}
\newtheorem{corollary}[theorem]{Corollary}
\newcommand{\F}{\mathbb F}
\newcommand{\mr}{\operatorname{mr}}
\title{Affine quadratic bounds and signed transposition quotients of knot groups}
\author{Internal research working proof}
\date{16 November 2026}
\begin{document}
\maketitle

The quadratic-refinement bound has explicit sharp knot representations in
arbitrarily large dimension. The representations below have image the finite
Coxeter group $W(D_r)$ and apply to a subfamily of pretzel knots already
covered by established meridional-rank theorems. They turn the conditional
criterion into concrete cases, without claiming new cases of the conjecture.
An exact linear lifting test provides a way to seek further examples from
symmetric quotients. The full conjecture remains open. All uniform results
below rest on written proofs; the accompanying finite controls are not formal
verification or external review. Priority of the elementary reformulation
has not been established.

\section{The quadratic bound with a possibly degenerate form}

Let $V=\F_2^k$, and let $B$ be a fixed \emph{nonzero} alternating bilinear
form on $V$. Nondegeneracy is not required. Suppose a group $H$ acts
transitively on $V$ and has generators
\[
 T_i(x)=x+(B(u_i,x)+a_i)u_i,
 \qquad 0\ne u_i\in V,\quad a_i\in\F_2.                 \tag{1}
\]
When $B(u_i,-)\ne0$, these are affine transvections. Every linear part
preserves $B$. The formula is stable under affine conjugation with
$B$-preserving linear part: conjugating by $h(x)=Lx+t$ replaces
$(u,a)$ by $(Lu,a+B(Lu,t))$.

\begin{theorem}\label{thm:quadratic}
Any generating family of the form (1) for a transitive $H$ has at least
$k+1$ members. In particular, a knot group with such a transitive
representation, with every meridian of the form (1), satisfies
$\mr(K)\ge k+1$, and $b(K)\ge k+1$.
\end{theorem}
\begin{proof}
Each generator moves points only within cosets of
$U=\langle u_1,\ldots,u_w\rangle$. Transitivity forces $U=V$, hence
$w\ge k$. If $w=k$, these directions are a basis. Every alternating form,
including a degenerate one, has quadratic refinements $q$ with $q(0)=0$;
their values on a basis can be prescribed arbitrarily. Prescribe
$q(u_i)=1+a_i$. Writing $b=B(u_i,x)$ gives
\[
 q(T_ix)+q(x)=(b+a_i)(q(u_i)+b)
             =(b+a_i)(1+a_i+b)=0.
\]
Thus $H$ preserves $q$. The polar form of a constant normalized $q$ is zero,
whereas $B\ne0$. Consequently $q$ has two nonempty level sets, contradicting
transitivity. Apply this to images of any meridional generating family.
The standard bridge presentation gives $\mr(K)\le b(K)$.
\end{proof}

This extends the nondegenerate argument in
\href{../cross-programme-connections/BREAKTHROUGH_SEARCH_TOPOLOGICAL_COUPLING_2026_11_16.tex}{the original coupling proof}.
The condition $B\ne0$ is necessary: for $B=0$, $k$ independent translations
act transitively. Inconsistency of the refinement equations only rules out
this particular obstruction; it does not prove transitivity for arbitrary
affine transvection groups.

\section{Signed transpositions and the exact transitivity criterion}

For $r\ge3$, set
\[
 V_r=\{x\in\F_2^r:\textstyle\sum_jx_j=0\},\quad
 B(x,y)=\textstyle\sum_jx_jy_j,\quad k=r-1.
\]
On $V_r$, $B$ is alternating and nonzero. Its radical is zero for odd $r$,
and $\langle(1,\ldots,1)\rangle$ for even $r$.
For $p\ne q$ put $u_{pq}=e_p+e_q$ and
\[
 s_{pq}^a(x)=x+(x_p+x_q+a)(e_p+e_q)
            =(pq)x+a(e_p+e_q),\qquad a\in\F_2.          \tag{2}
\]
These are involutions preserving $V_r$. They are of the form (1) and have
$2^{r-2}$ fixed points. They lie in
\[
 D_r=V_r\rtimes S_r=W(D_r),\qquad |D_r|=2^{r-1}r!.
\]
For $r=3$, $W(D_3)\cong W(A_3)\cong S_4$.
The action is faithful: a coordinate permutation fixing every $e_p+e_q$
is the identity when $r\ge3$, and a translation fixing zero is zero.

Make a multigraph on $\{1,\ldots,r\}$ from generators $s_{p_iq_i}^{a_i}$,
allowing repeated edges. Call its signing \emph{balanced} if there is
$b\in\F_2^r$ with $a_i=b_{p_i}+b_{q_i}$ on every edge. Equivalently every
cycle, including a length-two parallel-edge cycle, has sign sum zero.
The underlying graph is connected exactly when the projected transpositions
generate $S_r$.

\begin{theorem}\label{thm:signedgraph}
For a connected underlying graph and $r\ge3$, the following are equivalent:
the signing is unbalanced; the generated group is all of $D_r$; the action
on $V_r$ is transitive.
\end{theorem}
\begin{proof}
Choose a spanning tree. An ambient translation $x\mapsto x+b$, with
$b\in\F_2^r$, changes edge signs to $a_i+b_{p_i}+b_{q_i}$.
Choose $b$ to make every tree sign zero. The resulting tree swaps generate
the entire linear $S_r$. If the signing is unbalanced there is a remaining
sign-one edge. Multiply its signed swap by the corresponding linear swap
to obtain translation by $e_p+e_q$. Conjugation by $S_r$ supplies all
root translations; they span $V_r$. Thus the translated group is $D_r$.
Ambient translations normalize $D_r$ (since $b+\sigma b\in V_r$), so the
original group is also $D_r$ and is transitive.

If the signing is balanced, the original group preserves the Hamming
weight of $x+b$. As $x$ runs through $V_r$, $x+b$ runs through a parity
coset. An even coset has weights zero and two; an odd coset has weights
one and three. For $r\ge3$, each has at least two weights, precluding
transitivity. This also justifies the argument when $b$ is odd and hence
is not an origin shift inside $V_r$.
\end{proof}

\begin{corollary}
If a knot group surjects onto $D_r$ with meridians signed transpositions,
then $\mr(K)\ge r$. A certified upper bound $b(K)\le r$ proves equality.
For odd $r$ this uses the original nondegenerate theorem; even $r$ uses
Theorem~\ref{thm:quadratic} with one-dimensional radical.
\end{corollary}

This is also a Coxeter quotient bound. The usual real reflection
representation of $W(D_r)$ gives the same lower bound directly: fewer
than $r$ reflecting hyperplanes have a nonzero common fixed vector.
Thus this family of affine representations implements a known quotient
method, rather than bypassing its scope limitations.

\section{Explicit sharp pretzel representations in every rank}

\begin{theorem}\label{thm:pretzel}
Let $r\ge3$, $m\in\mathbb Z\setminus\{0\}$ and
$h_1,\ldots,h_{r-1}\in\mathbb Z$. The pretzel knot
\[
 K=P\bigl(2m,\,3(2h_1+1),\ldots,3(2h_{r-1}+1)\bigr)    \tag{3}
\]
admits a transitive signed-transposition representation on $2^{r-1}$
sheets. Its meridional rank and bridge number are both $r$.
\end{theorem}
\begin{proof}
There is exactly one even pretzel parameter; following the strand
connections therefore gives one component. All parameters have absolute
value at least two. Section 2 of Baader--Blair--Kjuchukova
\cite{bbk} supplies a meridian-preserving epimorphism onto the cycle group
\[
 C=\langle t_0,\ldots,t_{r-1}\mid t_i^2=1,
       (t_it_{i+1})^{|p_i|}=1\ (i\bmod r)\rangle.       \tag{4}
\]
Pairs not consecutive have no imposed relation in (4).
In particular, one must not add commutation relations on those pairs.
The same pretzel diagram has $r$ maxima, giving $b(K)\le r$.

Use coordinates numbered $0,\ldots,r-1$ and define
\[
 t_0\longmapsto s_{0,r-1}^0,\quad
 t_1\longmapsto s_{0,r-1}^1,\quad
 t_j\longmapsto s_{j-1,r-1}^0\quad(2\le j\le r-1).
                                                               \tag{5}
\]
The first pair multiplies to a root translation of order two. Every other
consecutive pair uses distinct roots sharing the last coordinate and
has product of order three. To check the latter, any two such signed
swaps can be made ordinary swaps by an ambient translation, so their
product is conjugate to a three-cycle. Thus all relations (4) hold for
the parameters (3), including negative twists.

The multigraph in (5) is connected, and its repeated edge has opposite
signs. It is unbalanced. Theorem~\ref{thm:signedgraph} proves that the
image is $D_r$, with a transitive action on $2^{r-1}$ sheets.
Theorem~\ref{thm:quadratic} gives $\mr(K)\ge r$, matching the upper bound.
\end{proof}

For example, $P(2,3,3,3,3)$ has an explicit transitive 16-sheet
affine-symplectic representation proving $\mr=b=5$;
$P(2,3,3,3,3,3,3)$ has a 64-sheet example proving $\mr=b=7$.
The even-rank example $P(2,3,3,3)$ uses eight sheets and the degenerate
extension, proving $\mr=b=4$.
These equalities already follow from the pretzel and Montesinos literature;
the contribution here is explicit realization of the proposed affine bound.

\section{A linear test that raises a symmetric quotient bound by one}

Suppose a presentation $G=\langle m_1,\ldots,m_s\mid R\rangle$ is supplied,
every generator is a meridian, and a specified quotient
$\phi:G\twoheadrightarrow S_r$ sends $m_i$ to $(p_iq_i)$.
Seek images $m_i\mapsto s_{p_iq_i}^{a_i}$.
Multiplication in $V_r\rtimes S_r$ is
$(z,P)(w,Q)=(z+Pw,PQ)$. Evaluating relators therefore imposes homogeneous
linear equations $Ma=0$ over $\F_2$. The permutation part of every
relator must be checked first. Negative letters cause no ambiguity since
all proposed generator images are involutions.

Let $Z=\ker M\subset\F_2^s$, and let
\[
 C_\phi=\{(b_{p_i}+b_{q_i})_{i=1}^s:b\in\F_2^r\}.
\]
Every ambient cut vector lies in $Z$ by affine conjugation of the linear
quotient. Connectedness gives $\dim C_\phi=r-1$.
\begin{corollary}\label{cor:algorithm}
The specified symmetric quotient admits a transitive signed-transposition
lift if and only if $\dim Z>r-1$. Any vector in $Z\setminus C_\phi$ is
a certificate. For a knot group such a certificate raises the usual
transposition quotient bound $\mr\ge r-1$ to $\mr\ge r$.
\end{corollary}
\begin{proof}
The relator equations are necessary and sufficient for a homomorphism
from the supplied presentation. Theorem~\ref{thm:signedgraph} says exactly
which homomorphisms are transitive.
\end{proof}

For even $r$, quotienting only by coboundaries from $b\in V_r$ is
incorrect. An odd ambient potential can produce a balanced action with
no common fixed point in $V_r$ but several orbits. For $r=4$, the
generators $s_{03}^1,s_{13}^0,s_{23}^0$ have no common fixed point on eight
sheets, yet have two four-point orbits. The correct cut space uses the
whole ambient $\F_2^r$.

\section{A uniform obstruction to the entire affine criterion}

\begin{lemma}\label{lem:torus}
Let $G_{p,q}=\langle x,y\mid x^p=y^q\rangle$ be a torus knot group and
let $\rho:G_{p,q}\to H$ have finite image. If $\gcd(q,|H|)=1$, the image
is abelian. If moreover meridians map to involutions, its order is at
most two.
\end{lemma}
\begin{proof}
The element $z=x^p=y^q$ is central. Since $q$ is relatively prime to
the order of $\rho(y)$, the latter is a power of $\rho(y)^q=\rho(z)$
and is therefore central. The image is generated by $\rho(x)$ and the
central $\rho(y)$, so is abelian. A knot group has abelianization $\mathbb Z$
generated by a meridian; an abelian image generated by a meridian
involution has order at most two.
\end{proof}

\begin{theorem}\label{thm:nouniversal}
For every $p\ge3$ and every prime $q>2^{p-1}-1$, the knot $T(p,q)$
has no transitive representation satisfying Theorem~\ref{thm:quadratic}
in any dimension. In particular, neither the original nondegenerate
criterion nor its degenerate extension is universal for knots.
\end{theorem}
\begin{proof}
The prime $q$ is greater than $p$, so these are coprime torus parameters.
The usual $p$-strand diagram gives $b(T(p,q))\le p$. If a representation
in dimension $k$ existed, Theorem~\ref{thm:quadratic} would give
$k+1\le\mr\le p$, hence $2\le k\le p-1$ (a nonzero alternating form
requires dimension at least two). But
\[
 |\operatorname{AGL}(k,2)|
   =2^{k+k(k-1)/2}\prod_{j=1}^{k}(2^j-1).
\]
Every odd prime divisor is at most $2^k-1\le2^{p-1}-1<q$.
The image order therefore is coprime to $q$.
Every map (1) is an involution because $B(u_i,u_i)=0$.
Lemma~\ref{lem:torus} forces image order at most two, which cannot
act transitively on $2^k\ge4$ points.
\end{proof}

For example $T(3,5)$ admits no representation of the stipulated type
in any dimension, though its known meridional rank and bridge number
are both three. This obstruction applies to arbitrary affine images,
not only Coxeter images. It is a limitation of the proposed criterion,
not a counterexample to MRC. The elementary group-order deduction is
recorded without a novelty claim.

\section{Verification and the remaining conjecture}

\texttt{signed\_lifts.py} implements Corollary~\ref{cor:algorithm} using
standard-library integer bit arithmetic. Its input must carry independent
provenance for the group presentation and the meridian labels. The saved
pretzel example is a presentation of the published \emph{quotient} (4),
not a full presentation of a knot group; composition with the published
epimorphism is the reason its certificate applies to the knot.

\texttt{verify\_signed\_lifts.py} exhausts connected simple signed graphs
on three and four vertices, checks repeated-edge graphs and the pretzel
quotients in ranks 3 through 14, independently enumerates the groups in
ranks 3 through 5, and retains negative controls. Changing the first
twist from two to three forces the two parallel signs to agree and
removes every transitive lift of that \emph{fixed} symmetric quotient.
This is not nonexistence of all possible quotients of that knot.

No arbitrary knot is proved to possess a sharp representation. The next
concrete gate is to obtain certified symmetric transposition colourings
and matching bridge upper bounds for knots outside existing resolved
families, then run the linear test. If it fails, other symmetric
colourings, other quotient groups, or an entirely different lower-bound
method may still succeed. Finite Coxeter quotients are not universal:
known torus knots with no nontrivial Coxeter quotient already rule out
a proof of the full conjecture restricted to the $W(D_r)$ construction
\cite{bkm}. Theorem~\ref{thm:nouniversal} gives an independent obstruction
to universality even after admitting arbitrary affine-symplectic images.

\begin{thebibliography}{9}
\bibitem{bbk} S. Baader, R. Blair, A. Kjuchukova,
Coxeter groups and meridional rank of links,
Math. Ann. 379 (2021), 1533--1551, especially Section 2.
\url{https://doi.org/10.1007/s00208-020-02124-z}.
\bibitem{bkm} R. Blair, A. Kjuchukova, N. Morrison,
Coxeter quotients of knot groups through 16 crossings,
especially Section 5 and Theorem 5.3.
\url{https://arxiv.org/abs/2208.09032}.
\end{thebibliography}
\end{document}
